% !PRINTED 559
\weberpage{559}

%<<CONT
\[
\begin{aligned}
(a_{0,0} - \eta)\,\eta + &\ a_{1,0}\eta_1 + \cdots \hspace{2em} a_{e-1,0}\eta_{e-1} = 0 \\
a_{0,1}\eta \ + &\ (a_{1,1} - \eta)\eta_1 + \cdots \hspace{2em} a_{e-1,1}\eta_{e-1} = 0 \\
&\cdots\cdots\cdots\cdots\cdots\cdots\cdots\cdots\cdots\cdots \\
a_{0,e-1} \ + &\ a_{1,e-1}\eta_1 + \cdots (a_{e-1,e-1} - \eta)\eta_{e-1} = 0,
\end{aligned}
\]
% [sic] The printed last row omits the factor \eta after a_{0,e-1}; retained diplomatically although the parallel rows require it.
so sieht man, da die $\eta, \eta_1 \ldots \eta_{e-1}$ nicht verschwinden, dass die Determinante dieses Systems Null sein muss, also

\weq{17}{\begin{vmatrix} a_{0,0}-\eta, & a_{1,0}, & \ldots & a_{e-1,0} \\ a_{0,1}, & a_{1,1}-\eta, & \ldots & a_{e-1,1} \\ \cdots & \cdots & \cdots & \cdots \\ a_{0,e-1}, & a_{1,e-1}, & \ldots & a_{e-1,e-1}-\eta \end{vmatrix} = 0,}

was, wenn man $x$ für $\eta$ schreibt, vom Vorzeichen abgesehen, mit $F_e(x)$ übereinstimmen muss.

\gesperrt{Die Function $F_e(x)$ hat also nicht bloss rationale, sondern ganzzahlige Coëfficienten.}

% [sic] authority says that r is indeterminate here, although the displayed polynomial has indeterminate x and r is one of its roots.
Man kann nun auch die Gleichung $f^{\mathrm{ten}}$ Grades bilden, deren Wurzeln $r, r_e, r_{2e} \ldots r_{(f-1)e}$ sind; denn es ist für ein unbestimmtes $r$

\weq{18}{\Phi_e(x) = (x - r)(x - r_e) \cdots (x - r_{(f-1)e})}

eine Function im Körper $R(\eta)$. Ihre Coëfficienten lassen sich leicht aus den \gesperrt{Newton'schen} Formeln berechnen, durch die die Coëfficienten einer Gleichung mittelst der Potenzsummen der Wurzeln ausgedrückt werden (\S.\,42). Denn es ist nach (12) und (3)
\[
r^{g^h} + r_e^{g^h} + \cdots + r_{(f-1)e}^{g^h} = \eta_h,
\]
% [resolved 2026-08-06, 5x zoom strip 4/5] r^{g^h} is correct. The exponent really is g
% carrying a further superscript h, legible on all four terms of the display, and it agrees
% with the sense: the sum is eta_h, the h-th period, indexed by the same h. Stands.
also geradezu einer der Potenzsummen gleich, und jede Potenzsumme ist einer der Perioden gleich.

Die Gleichung (18) lässt sich aber ebenso behandeln wie die Gleichung (2); denn die Gruppe von $\Phi_e(x) = 0$ ist eben $C_e$, und wenn $e'$ ein Theiler von $f$ ist, etwa $f = e'f'$, so ist $C_{ee'}$ ein Theiler von $C_e$. Zu $C_{ee'}$ gehört aber die Periode
\[
\eta' = r + r_{ee'} + r_{2ee'} + \cdots + r_{(f'-1)ee'},
\]
die also von einer Gleichung des Grades $e'$ im Körper $R(\eta)$ abhängt.

Wenn man die Zahlen $e, e' \ldots$ als Primzahlen annimmt, so besteht die Reihe der Resolventen aus einer Reihe von Gleichungen, deren Grade die Primfactoren von $n - 1$ sind.

Was die Gruppe der Gleichung $e^{\mathrm{ten}}$ Grades betrifft, deren Wurzeln die $e$ Grössen $\eta$ sind, so finden wir diese sehr einfach
%CONT>>
